\documentclass[12pt,a4paper]{report}

% ============================================================
% Paquetes
% ============================================================
\usepackage[utf8]{inputenc}
\usepackage[T1]{fontenc}
\usepackage[spanish,es-noshorthands]{babel}
\usepackage[a4paper,margin=2.5cm]{geometry}
\usepackage{graphicx}
\usepackage{booktabs}
\usepackage{tabularx}
\usepackage{longtable}
\usepackage{listings}
\usepackage{xcolor}
\usepackage{hyperref}
\usepackage{enumitem}
\usepackage{fancyhdr}
\usepackage{titlesec}
\usepackage{parskip}

% ============================================================
% Datos del documento (editar aquí)
% ============================================================
\newcommand{\nombreEstudiante}{[Nombre del estudiante]}
\newcommand{\nombreDocente}{[Nombre del docente]}
\newcommand{\nombreAsignatura}{Desarrollo Fullstack I (DSY1103)}
\newcommand{\nombreInstitucion}{Duoc UC}
\newcommand{\fechaInforme}{\today}

% ============================================================
% Configuración visual
% ============================================================
\hypersetup{
    colorlinks=true,
    linkcolor=blue!50!black,
    citecolor=black,
    urlcolor=blue!50!black,
    pdftitle={BookPoint Chile - Informe Tecnico ET},
    pdfauthor={\nombreEstudiante}
}

\lstdefinelanguage{json}{
    basicstyle=\ttfamily\small,
    showstringspaces=false,
    breaklines=true,
    numbers=left,
    numberstyle=\tiny\color{gray},
    string=[s]{"}{"},
    stringstyle=\color{green!40!black},
    literate=
        *{0}{{{\color{blue!70!black}0}}}{1}
         {1}{{{\color{blue!70!black}1}}}{1}
         {2}{{{\color{blue!70!black}2}}}{1}
         {3}{{{\color{blue!70!black}3}}}{1}
         {4}{{{\color{blue!70!black}4}}}{1}
         {5}{{{\color{blue!70!black}5}}}{1}
         {6}{{{\color{blue!70!black}6}}}{1}
         {7}{{{\color{blue!70!black}7}}}{1}
         {8}{{{\color{blue!70!black}8}}}{1}
         {9}{{{\color{blue!70!black}9}}}{1},
}

\lstset{
    basicstyle=\ttfamily\small,
    breaklines=true,
    frame=single,
    columns=fullflexible,
    keywordstyle=\color{blue!70!black}\bfseries,
    commentstyle=\color{gray},
    stringstyle=\color{green!40!black},
    showstringspaces=false,
    numbers=left,
    numberstyle=\tiny\color{gray},
    tabsize=2
}

\pagestyle{fancy}
\fancyhf{}
\fancyhead[L]{BookPoint Chile - Informe Tecnico}
\fancyhead[R]{\leftmark}
\fancyfoot[C]{\thepage}

\titleformat{\chapter}[hang]{\normalfont\huge\bfseries}{\thechapter.}{20pt}{\huge}

% ============================================================
% Documento
% ============================================================
\begin{document}

% ---------- Portada ----------
\begin{titlepage}
    \centering
    \vspace*{2cm}

    {\Large \nombreInstitucion \par}
    \vspace{0.3cm}
    {\large Escuela de Informática y Telecomunicaciones \par}
    \vspace{2cm}

    {\Huge \bfseries BookPoint Chile \par}
    \vspace{0.5cm}
    {\Large Informe Técnico -- Arquitectura de Microservicios \par}
    \vspace{0.3cm}
    {\large Evaluación Final Transversal -- Situación Evaluativa 1: Entrega por Encargo \par}

    \vspace{3cm}

    {\large \nombreAsignatura \par}
    \vspace{0.5cm}
    {\large Estudiante: \nombreEstudiante \par}
    \vspace{0.3cm}
    {\large Docente: \nombreDocente \par}

    \vfill

    {\large \fechaInforme \par}
\end{titlepage}

\tableofcontents
\clearpage

% ---------- Secciones ----------
\chapter{Introducción}

\section{Contexto}

BookPoint Chile es una empresa nacional dedicada a la venta de libros, artículos de
papelería y material educativo. Nació como una librería independiente en Concepción y,
debido al crecimiento sostenido de sus ventas, abrió nuevas sucursales en Temuco y La
Serena.

Actualmente, la empresa opera con un sistema monolítico que concentra ventas,
inventario, despacho y atención de clientes en un único aplicativo. El aumento
sostenido de pedidos, tanto presenciales como por canal web, ha comenzado a provocar
lentitud en los procesos, problemas de stock desactualizado entre sucursales y demoras
en la gestión de despachos.

\section{Objetivo del informe}

Este informe documenta la propuesta de solución tecnológica desarrollada para
modernizar el sistema de BookPoint Chile, migrándolo desde su arquitectura monolítica
actual hacia una arquitectura de microservicios independientes. En concreto, cubre los
siete aspectos exigidos por la pauta de la Evaluación Final Transversal para la
Situación Evaluativa 1 (Entrega por Encargo):

\begin{enumerate}
    \item La estrategia de microservicios utilizada para resolver el caso propuesto.
    \item La evaluación de los enfoques éticos considerados en su implementación.
    \item Las buenas prácticas de diseño, arquitectura, y las herramientas (Spring,
    Maven, Git) utilizadas para desarrollar los componentes de backend.
    \item La implementación de operaciones CRUD con JPA/ORM, validadas con Postman.
    \item La comunicación RESTful entre microservicios, validada con Postman.
    \item El desarrollo de pruebas unitarias con JUnit.
    \item La implementación de un framework de documentación HATEOAS.
\end{enumerate}

\section{Alcance del proyecto}

La solución desarrollada consta de \textbf{diez microservicios de negocio}
(catálogo, inventario, sucursales, clientes, carrito, pedidos, pagos, envíos,
notificaciones y reseñas), además de un servidor de descubrimiento (Eureka) y un API
Gateway como punto de entrada único.

Se priorizó cubrir en profundidad el \textbf{flujo comercial completo} descrito en el
caso -- desde que un cliente explora el catálogo hasta que recibe su pedido y puede
dejar una reseña --, con validaciones reales entre servicios, manejo de errores
consistente y pruebas automatizadas. Quedan deliberadamente fuera del alcance actual
funciones administrativas de soporte descritas en el caso (autenticación y roles,
reportes de venta, promociones, facturación electrónica, transferencias entre
sucursales), que se detallan junto con su justificación en la
sección~\ref{ch:arquitectura}.

\chapter{Estrategia y Arquitectura de Microservicios}
\label{ch:arquitectura}

\section{Análisis del sistema actual}

Según el caso de estudio, BookPoint Chile nació como una librería independiente en
Concepción y, debido al crecimiento de sus ventas, abrió nuevas sucursales en Temuco y
La Serena. Actualmente opera con un \textbf{sistema monolítico} que concentra en un
mismo aplicativo las ventas, el inventario, el despacho y la atención de clientes.

El propio caso identifica tres síntomas concretos de que ese monolito ya no escala con
el crecimiento de la operación:

\begin{itemize}
    \item \textbf{Lentitud en los procesos}, provocada por el aumento sostenido de
    pedidos tanto presenciales como por canal web.
    \item \textbf{Stock desactualizado} entre sucursales, al no existir una fuente de
    verdad de inventario que se sincronice en tiempo real entre Concepción, Temuco y
    La Serena.
    \item \textbf{Demoras en la gestión de despachos}, al no contar con un módulo
    dedicado de logística que pueda evolucionar independientemente del resto del
    sistema.
\end{itemize}

Estos síntomas son característicos de una arquitectura monolítica sometida a
crecimiento: al compartir un único proceso, una única base de datos y un único ciclo de
despliegue, cualquier cambio en el módulo de inventario obliga a redesplegar
\textit{todo} el sistema -- incluida la caja de venta presencial, que no puede
detenerse sin afectar directamente la operación de las tres tiendas físicas. Del mismo
modo, un pico de tráfico en el canal web (por ejemplo, durante una campaña de
descuentos) satura recursos que también necesita la venta presencial, porque ambos
comparten el mismo proceso.

\section{Estrategia de microservicios elegida}

Para resolver el caso propuesto se optó por una arquitectura de
\textbf{microservicios independientes}, organizados por dominio de negocio (no por
capa técnica), con los siguientes principios:

\begin{itemize}
    \item \textbf{Base de datos por servicio.} Cada microservicio posee su propia base
    de datos MySQL (Tabla~\ref{tab:bd-servicio}), eliminando el acoplamiento de datos
    que hoy sufre el sistema monolítico. El módulo de inventario puede evolucionar su
    modelo de stock sin coordinar una migración con el módulo de reseñas.
    \item \textbf{Comunicación síncrona vía REST/Feign} para las relaciones entre
    dominios que sí necesitan datos de otro servicio en el momento (por ejemplo, un
    pedido necesita confirmar que el libro existe y que hay stock antes de
    concretarse).
    \item \textbf{Descubrimiento de servicios con Eureka}, para que ningún
    microservicio necesite conocer de antemano la dirección de red de los demás --
    relevante en un escenario con tres sucursales que podrían eventualmente desplegarse
    en ubicaciones distintas.
    \item \textbf{API Gateway como fachada única}, que permite exponer una sola puerta
    de entrada hacia los distintos canales (presencial, web) sin que el cliente
    necesite conocer la topología interna de microservicios.
\end{itemize}

Esta estrategia ataca directamente los tres síntomas identificados: la lentitud se
resuelve porque cada dominio escala y se despliega de forma independiente (el catálogo,
que se consulta constantemente, no compite por recursos con el módulo de pagos); el
stock desactualizado se resuelve centralizando el inventario en un único microservicio
consultado por todos los canales y sucursales; y las demoras de despacho se resuelven
con un microservicio de envíos dedicado, capaz de evolucionar su lógica de logística sin
tocar el resto del sistema.

\section{Diseño de los microservicios}

El caso de estudio describe seis perfiles de usuario (Administrador del Sistema, Jefe
de Sucursal, Asistente de Ventas, Área de Logística, Encargado de Bodega y Clientes vía
web) con un total de más de treinta acciones posibles. La Tabla~\ref{tab:mapeo-perfiles}
mapea los diez microservicios efectivamente construidos a los perfiles del caso que
cubren.

\begin{table}[h]
\centering
\small
\begin{tabularx}{\textwidth}{lX}
\toprule
\textbf{Microservicio} & \textbf{Perfiles/acciones del caso que cubre} \\
\midrule
Catálogo & Cliente: navegar catálogo de libros. \\
Inventario & Encargado de Bodega: administrar stock central y por sucursal.
Administrador: consultar stock disponible en otras sucursales. \\
Sucursales & Jefe de Sucursal: administrar datos de la sucursal. \\
Clientes & Cliente: crear cuenta, administrar perfil y direcciones. \\
Carrito & Cliente: agregar productos al carrito. \\
Pedidos & Cliente: realizar compras online. Asistente de Ventas: registrar ventas
(equivalente digital de la venta en caja). Jefe de Sucursal: generar reportes de
ventas por sucursal y por autor. \\
Pagos & Cliente: pagar una compra online. \\
Envíos & Área de Logística: gestionar envíos de pedidos a domicilio, actualizar
estados de pedidos. \\
Notificaciones & Cliente: recibir confirmaciones y alertas asociadas a su pedido. \\
Reseñas & Cliente: dejar reseñas y calificaciones de productos. \\
\bottomrule
\end{tabularx}
\caption{Mapeo de los microservicios construidos a los perfiles del caso BookPoint Chile}
\label{tab:mapeo-perfiles}
\end{table}

\subsection{Alcance y exclusiones deliberadas}

No todas las acciones descritas en el caso están cubiertas por la implementación
actual. Se priorizó el \textbf{flujo comercial completo} (explorar catálogo, agregar al
carrito, comprar, pagar, despachar, reseñar) por sobre funciones administrativas de
soporte, que quedaron fuera de alcance de forma consciente:

\begin{itemize}
    \item Gestión de usuarios y permisos por rol para el perfil Administrador del
    Sistema (autenticación básica de clientes vía JWT sí se implementó, ver
    sección~\ref{ch:etica}).
    \item Reportes de ventas por categoría (Jefe de Sucursal) -- los reportes por
    sucursal y por autor sí se implementaron (sección~\ref{ch:crud}), pero por
    categoría exigiría antes agregar un campo de categoría/género al catálogo, que
    no existe hoy.
    \item Promociones, descuentos, convenios estudiantiles y cupones (Asistente de
    Ventas, Cliente).
    \item Emisión de boletas o facturas electrónicas y gestión de devoluciones
    (Asistente de Ventas).
    \item Transferencias de stock entre sucursales, gestión de proveedores editoriales
    y optimización de rutas de distribución (Área de Logística, Encargado de Bodega).
\end{itemize}

Esta decisión de alcance no es un descuido: dado el tiempo disponible, se priorizó
profundizar el flujo transaccional principal -- con validaciones cruzadas reales entre
servicios, manejo de errores consistente y pruebas automatizadas -- por sobre cubrir
superficialmente la totalidad de acciones descritas en el caso.

\section{Plan de migración}

Para migrar el sistema monolítico descrito en el caso hacia esta arquitectura, se
plantea un enfoque incremental de tipo \textbf{strangler fig} (extracción progresiva),
en lugar de una reescritura total (\textit{big-bang rewrite}), que sería
significativamente más riesgosa para una empresa que necesita mantener sus tres
sucursales operativas durante la transición:

\begin{enumerate}
    \item \textbf{Extraer primero los dominios de solo lectura y bajo acoplamiento
    transaccional}: catálogo e inventario. Son los más consultados (cada venta,
    presencial o web, necesita leerlos) y los que menos dependen de otros módulos para
    funcionar.
    \item \textbf{Introducir el API Gateway como fachada de transición}, de modo que
    el resto del sistema (aún monolítico en esta etapa) y los microservicios ya
    extraídos convivan detrás de una única puerta de entrada, sin que el cliente note
    el cambio.
    \item \textbf{Extraer los dominios transaccionales}: clientes, carrito, pedidos y
    pagos, en ese orden, ya que cada uno depende de que el anterior esté disponible
    como servicio (el carrito necesita clientes y catálogo; los pedidos necesitan
    catálogo e inventario).
    \item \textbf{Extraer los dominios de soporte}: sucursales, envíos, notificaciones
    y reseñas, que dependen de que pedidos ya exista como microservicio independiente.
    \item \textbf{Retirar el monolito} una vez que todos los dominios estén servidos
    por microservicios y validados en producción.
\end{enumerate}

Este orden minimiza el riesgo operativo: en cada paso, el sistema sigue siendo
funcional para las tres sucursales, y solo se apaga definitivamente el monolito cuando
ya no queda ningún flujo de negocio dependiendo de él.

\chapter{Evaluación de Enfoques Éticos}
\label{ch:etica}

\section{Privacidad de datos de clientes}

El microservicio de clientes almacena datos personales identificables: nombre, correo
electrónico, dirección y comuna. Al tratarse de un sistema que atiende tanto canal
presencial como online, estos datos circulan entre varios microservicios (por ejemplo,
\texttt{envíos} necesita la dirección de destino asociada a un pedido).

Frente a esto se aplicó un principio de \textbf{minimización de datos}: cada
microservicio conoce únicamente los campos estrictamente necesarios para su función. El
microservicio de pagos, por ejemplo, no tiene acceso al nombre ni a la dirección del
cliente -- solo al identificador del pedido y al monto a cobrar --, de modo que un
eventual incidente de seguridad en el módulo de pagos no expondría por sí solo datos
personales identificables.

Inicialmente el sistema no implementaba autenticación ni cifrado de credenciales,
decisión de alcance consciente dado el tiempo disponible del proyecto y no un descuido.
Sin embargo, precisamente por tratarse de datos personales de clientes -- direcciones,
historial de compras --, se decidió cerrar esa brecha al menos en el punto donde más
importa: el microservicio de clientes ahora exige un token JWT válido para modificar o
eliminar un perfil, y las contraseñas se almacenan con \texttt{BCryptPasswordEncoder}
(hash con salt, nunca en texto plano). La protección se limitó deliberadamente a las
operaciones de escritura sobre el propio perfil (\texttt{PUT}/\texttt{DELETE}) y no se
extendió a las lecturas ni al resto de los microservicios: \texttt{carrito} y
\texttt{reseñas} consultan clientes vía Feign para sus propias validaciones cruzadas, y
esas llamadas de servicio a servicio no llevan el token de ningún usuario. Proteger
también las lecturas habría exigido propagar un token de servicio a través de Feign --
una solución más completa, pero fuera del alcance de tiempo disponible. Esta es, en sí
misma, una decisión ética explícita: priorizar cerrar la brecha más sensible (que
cualquiera pueda modificar los datos de otro cliente) por sobre una implementación
completa pero más costosa.

\section{Responsabilidad en el despliegue}

BookPoint Chile opera simultáneamente en tres sucursales físicas y un canal web. A
diferencia de un negocio puramente digital, una caída del sistema no solo afecta las
ventas online: si el monolito actual falla, se detiene también la venta presencial en
Concepción, Temuco y La Serena, porque todo depende del mismo proceso.

La arquitectura de microservicios propuesta mitiga parcialmente este riesgo mediante
\textbf{aislamiento de fallas}: si el microservicio de reseñas deja de responder, el
cliente puede seguir explorando el catálogo, agregando productos al carrito y pagando
su compra sin verse afectado. Esto traslada la responsabilidad del despliegue desde "no
puede fallar nunca" (inviable en cualquier sistema real) hacia "una falla puntual no
debe paralizar el negocio completo" -- un estándar más realista y, a la vez, más
exigente en el diseño de cada servicio individual, porque cada uno debe manejar
adecuadamente los casos en que sus dependencias no respondan (ver
sección~\ref{ch:rest}, manejo de errores en comunicación remota).

\section{Impacto en el empleo}

El sistema automatiza tareas hoy realizadas manualmente por el personal de BookPoint:
el registro de ventas en caja (equivalente al microservicio de pedidos), la consulta de
disponibilidad de stock entre sucursales (inventario) y el seguimiento del estado de un
envío (envíos, con su historial de estados).

Es importante encuadrar este punto con honestidad: el sistema no reemplaza al
Asistente de Ventas, al Encargado de Bodega ni al Área de Logística -- automatiza
tareas repetitivas y propensas a error humano (como mantener sincronizado el stock
entre tres sucursales a mano), liberando tiempo para tareas que sí requieren juicio
humano: atención al cliente, resolución de casos excepcionales, gestión de proveedores.
El riesgo ético real no está en automatizar estas tareas puntuales, sino en cómo se
gestione la transición del personal hacia esas nuevas responsabilidades -- una decisión
organizacional de BookPoint Chile que excede el alcance técnico de este informe, pero
que vale la pena nombrar explícitamente como parte de la evaluación ética del proyecto.

\section{Desafíos éticos específicos enfrentados durante el desarrollo}

Más allá de los principios generales discutidos arriba, el desarrollo del proyecto
presentó tensiones concretas entre avanzar rápido y hacerlo bien, y al menos una duda
genuina sobre si una decisión de alcance era aceptable para un sistema real.

\textbf{Errores que exigieron parar y corregir en vez de seguir avanzando.} En varios
momentos del desarrollo surgieron problemas que hubiera sido más rápido ignorar o
parchar superficialmente, pero que se decidió corregir de raíz:

\begin{itemize}
    \item Durante la mayor parte del desarrollo, cualquier recurso no encontrado
    devolvía \texttt{500} en lugar de \texttt{404} en los diez servicios. Era un error
    silencioso -- el sistema "funcionaba" en el sentido de que no se caía --, pero
    devolver el código HTTP correcto es justamente lo que permite que un cliente de la
    API distinga un error suyo de un error del servidor. Se optó por parar y corregirlo
    en los diez servicios a la vez, en lugar de dejarlo para "después".
    \item Un cliente de la API podía enviar un \texttt{id} propio al crear un recurso, y
    ese \texttt{id} terminaba sobrescribiendo silenciosamente un recurso existente en
    vez de crear uno nuevo. Es el tipo de bug que puede pasar desapercibido en las
    demos (todo parece funcionar) pero que en un sistema real significaría corromper
    datos de otro cliente sin ningún aviso.
    \item Al centralizar la documentación Swagger de los diez microservicios a través
    del Gateway, las rutas no llegaban correctamente: hubo que re-enrutar manualmente
    los puertos en el YAML del Gateway y ajustar la configuración de cada
    microservicio para que calzara con lo que el Gateway esperaba. Y al containerizar
    los servicios, Swagger dejó de funcionar porque cada servicio reportaba su propia
    IP interna de Docker como servidor, en vez de una URL que el navegador del cliente
    pudiera alcanzar. Ninguno de los dos problemas era obvio a primera vista, y ambos
    se podrían haber dejado "así no más" con la excusa de que la API seguía
    respondiendo -- se optó por diagnosticar la causa real en vez de silenciar el
    síntoma.
    \item La decisión de alcance de la autenticación JWT (sección~\ref{ch:etica}) es en
    sí misma un ejemplo de esta tensión: la alternativa más rápida habría sido proteger
    todo el microservicio de clientes sin pensar en las consecuencias; se optó por
    invertir tiempo en entender qué llamadas dependían de un acceso sin token
    (\texttt{carrito} y \texttt{reseñas} vía Feign) antes de decidir el alcance
    correcto de la protección.
\end{itemize}

\textbf{Duda genuina sobre decisiones de alcance.} La ausencia inicial de
autenticación fue la que generó más incomodidad real durante el desarrollo: saber que,
mientras no se implementara, cualquiera con acceso a la API podía consultar o
modificar los datos de cualquier cliente sin restricción, se sentía como una brecha
real y no solo teórica -- razón por la cual se decidió cerrarla al menos parcialmente
antes de la entrega final, en lugar de dejarla únicamente documentada como limitación.

Una duda similar surgió con la exclusión de reportes de venta y de la gestión de roles
administrativos (sección~\ref{ch:arquitectura}): un Jefe de Sucursal o un Administrador
del Sistema reales necesitarían esas funciones para operar el negocio día a día, y su
ausencia no es solo una limitación técnica, sino una función de negocio real que el
caso pide. A diferencia de la autenticación, esta duda llevó a revisar el alcance
antes de la entrega final: los reportes de ventas por sucursal y por autor se
implementaron (sección~\ref{ch:crud}), por ser acotados y no requerir cambios de
arquitectura. La gestión de usuarios y permisos por rol, en cambio, se decidió
dejar fuera deliberadamente -- a diferencia de un reporte, que es una vista de solo
lectura sobre datos que ya existen, un sistema de roles reales implica un modelo de
autorización completo y una decisión arquitectónica (¿es un nuevo microservicio? ¿se
integra a uno existente?) que merece más tiempo de análisis del que quedaba
disponible, y apurarla habría sido peor que documentarla honestamente como
limitación.

\chapter{Buenas Prácticas de Desarrollo Backend}
\label{ch:buenas-practicas}

Este capítulo aborda el indicador de mayor ponderación del informe (16\%): el desarrollo
de componentes de backend aplicando buenas prácticas de diseño y arquitectura, el uso de
Spring y Maven, y el trabajo colaborativo mediante Git.

\section{Patrón Controller-Service-Repository}

Los diez microservicios de BookPoint aplican de forma consistente el patrón
\textbf{Controller-Service-Repository (CSR)}, separando estrictamente tres
responsabilidades:

\begin{itemize}
    \item \textbf{Controller}: recibe la petición HTTP, delega en el service y traduce
    el resultado a una respuesta REST (código de estado, cuerpo, enlaces HATEOAS). No
    contiene lógica de negocio.
    \item \textbf{Service}: concentra la lógica de negocio, las validaciones cruzadas
    contra otros microservicios (vía Feign) y las reglas del dominio (por ejemplo,
    verificar stock antes de confirmar un pedido).
    \item \textbf{Repository}: acceso a datos mediante Spring Data JPA
    (\texttt{JpaRepository}), sin lógica de negocio.
\end{itemize}

El siguiente ejemplo, tomado del microservicio de catálogo, muestra la capa de servicio
completa de \texttt{LibroServiceImpl}: no contiene código HTTP ni SQL, solo la
traducción entre el DTO de entrada y la entidad persistida, y el manejo del caso "no
encontrado" mediante una excepción de negocio propia:

\begin{lstlisting}[language=Java, caption={LibroServiceImpl.java -- capa de servicio de catálogo}]
@Service
@RequiredArgsConstructor
public class LibroServiceImpl implements LibroService {

    private final LibroRepository libroRepository;

    @Override
    public Libro guardarLibro(LibroDTO libroDTO) {
        Libro libro = new Libro();
        libro.setTitulo(libroDTO.getTitulo());
        libro.setAutor(libroDTO.getAutor());
        libro.setPrecio(libroDTO.getPrecio());
        return libroRepository.save(libro);
    }

    @Override
    public Libro obtenerPorId(Long id) {
        return libroRepository.findById(id)
            .orElseThrow(() -> new RecursoNoEncontradoException(
                "Libro no encontrado con id: " + id));
    }
}
\end{lstlisting}

Esta separación se replica de forma idéntica en los diez servicios, lo que permite que
cualquier persona familiarizada con uno de ellos pueda orientarse rápidamente en los
demás -- una consecuencia directa de mantener una convención uniforme en todo el
proyecto.

\section{Modularidad, mantenibilidad y escalabilidad}

\subsection{Base de datos por servicio}

Cada microservicio posee su propia base de datos MySQL, sin compartir esquemas ni
tablas con otros servicios (Tabla~\ref{tab:bd-servicio}). Esto evita el acoplamiento de
datos característico de los monolitos -- como el descrito en el caso BookPoint Chile,
donde ventas, inventario, despacho y atención de clientes conviven en un mismo sistema
-- y permite que cada equipo pueda evolucionar el modelo de datos de su dominio sin
coordinar migraciones con el resto.

\begin{table}[h]
\centering
\begin{tabular}{ll}
\toprule
\textbf{Microservicio} & \textbf{Base de datos} \\
\midrule
catálogo & bookpoint\_catalogo\_db \\
inventario & bookpoint\_inventario\_db \\
pedidos & bookpoint\_pedidos\_db \\
clientes & bookpoint\_clientes\_db \\
sucursales & bookpoint\_sucursales\_db \\
carrito & bookpoint\_carrito\_db \\
pagos & bookpoint\_pagos\_db \\
envíos & bookpoint\_envios\_db \\
notificaciones & bookpoint\_notificaciones\_db \\
reseñas & bookpoint\_resenas\_db \\
\bottomrule
\end{tabular}
\caption{Una base de datos independiente por microservicio}
\label{tab:bd-servicio}
\end{table}

\subsection{Separación DTO / entidad}

En los servicios donde la entrada requiere validación estricta (por ejemplo,
\texttt{catálogo}, \texttt{inventario} y \texttt{pedidos}), se utiliza un DTO anotado
con \textit{Bean Validation} (\texttt{@NotBlank}, \texttt{@NotNull}, \texttt{@Positive})
en lugar de exponer la entidad JPA directamente en el controlador. Esto evita que un
cliente de la API pueda, por ejemplo, enviar un \texttt{id} propio al crear un recurso
y sobrescribir accidentalmente uno existente -- un error real detectado y corregido
durante el desarrollo (ver sección~\ref{sec:evolucion}).

\subsection{Manejo centralizado de excepciones}

Cada microservicio expone un \texttt{@RestControllerAdvice} propio
(\texttt{GlobalExceptionHandler}) que traduce las excepciones de negocio en respuestas
HTTP coherentes, evitando repetir bloques \texttt{try/catch} en cada controlador:

\begin{lstlisting}[language=Java, caption={GlobalExceptionHandler.java (catálogo) -- fragmento}]
@RestControllerAdvice
public class GlobalExceptionHandler {

    @ExceptionHandler(RecursoNoEncontradoException.class)
    public ResponseEntity<ErrorResponseDTO> handleRecursoNoEncontrado(
            RecursoNoEncontradoException ex) {
        ErrorResponseDTO errorDTO = new ErrorResponseDTO(
                ex.getMessage(), HttpStatus.NOT_FOUND.value(),
                LocalDateTime.now(), null);
        return new ResponseEntity<>(errorDTO, HttpStatus.NOT_FOUND);
    }

    @ExceptionHandler(IllegalArgumentException.class)
    public ResponseEntity<ErrorResponseDTO> handleArgumentoInvalido(
            IllegalArgumentException ex) {
        ErrorResponseDTO errorDTO = new ErrorResponseDTO(
                ex.getMessage(), HttpStatus.BAD_REQUEST.value(),
                LocalDateTime.now(), null);
        return new ResponseEntity<>(errorDTO, HttpStatus.BAD_REQUEST);
    }
}
\end{lstlisting}

Este patrón asegura que un recurso inexistente responda siempre con \texttt{404}, una
violación de regla de negocio con \texttt{400}, y un error inesperado con \texttt{500}
-- de forma consistente en los diez servicios, en vez de depender de que cada
desarrollador recuerde hacerlo caso a caso.

\section{Spring Boot y Maven}

El proyecto está organizado como un \textbf{POM multi-módulo} de Maven: un
\texttt{pom.xml} raíz de tipo \texttt{pom} declara los doce módulos (diez
microservicios de negocio, \texttt{eureka-server} y \texttt{gateway}) y centraliza en
\texttt{dependencyManagement} las versiones compartidas -- \textbf{Spring Boot 4.1},
\textbf{Java 21} y \textbf{Spring Cloud 2025.1.2} -- de modo que ningún módulo hijo
necesita fijar versiones por su cuenta:

\begin{lstlisting}[language=XML, caption={pom.xml raíz -- fragmento relevante}]
<properties>
    <java.version>21</java.version>
    <spring-cloud.version>2025.1.2</spring-cloud.version>
</properties>

<dependencyManagement>
    <dependencies>
        <dependency>
            <groupId>org.springframework.cloud</groupId>
            <artifactId>spring-cloud-dependencies</artifactId>
            <version>${spring-cloud.version}</version>
            <type>pom</type>
            <scope>import</scope>
        </dependency>
    </dependencies>
</dependencyManagement>
\end{lstlisting}

Además de \texttt{spring-boot-starter-web} y \texttt{spring-boot-starter-data-jpa}
(presentes en los diez servicios), el proyecto utiliza:

\begin{itemize}
    \item \textbf{Spring Cloud OpenFeign}: comunicación declarativa entre
    microservicios (capítulo~\ref{ch:rest}).
    \item \textbf{Spring Cloud Netflix Eureka Client}: registro y descubrimiento de
    servicios.
    \item \textbf{Spring HATEOAS}: navegabilidad de la API (capítulo~\ref{ch:hateoas}).
    \item \textbf{springdoc-openapi}: documentación Swagger/OpenAPI por servicio.
    \item \textbf{JaCoCo Maven Plugin}: medición de cobertura de pruebas
    (capítulo~\ref{ch:testing}).
    \item \textbf{Datafaker}: generación de datos de prueba realistas al levantar cada
    servicio (ver sección~\ref{sec:evolucion}).
\end{itemize}

\section{Trabajo colaborativo con Git}

Este proyecto fue desarrollado íntegramente en solitario: aunque la pauta permite
equipos de hasta tres integrantes, todo el desarrollo backend -- diseño, código,
pruebas y correcciones -- fue trabajo individual, sin otros integrantes aportando
código al repositorio.

El trabajo se hizo directo sobre la rama \texttt{main}, sin ramas separadas por
funcionalidad. El desarrollo comenzó por los servicios núcleo del flujo comercial
(catálogo, inventario y pedidos), construidos hasta tener la comunicación Feign básica
funcionando entre ellos; a partir de ahí se fue extendiendo el resto de los
microservicios y, en tareas más mecánicas y repetitivas -- aplicar el mismo cambio de
forma consistente a lo largo de los diez servicios, como la incorporación de HATEOAS o
la construcción de la colección Postman --, se aceleró el trabajo con herramientas de
asistencia de desarrollo.

El repositorio del proyecto utiliza Git como sistema de control de versiones, con una
convención de mensajes de commit en español, explicando siempre el \textit{porqué} del
cambio y no solo el \textit{qué}. Los cambios no relacionados se separan en commits
distintos en lugar de agruparse en una sola entrega masiva -- por ejemplo, la sesión de
trabajo dedicada a corregir el mapeo de errores HTTP, agregar HATEOAS y sumar
validaciones cruzadas vía Feign se dividió en tres commits independientes, cada uno
enfocado en un solo cambio lógico:

\begin{lstlisting}[caption={Extracto de \texttt{git log --oneline}}]
47bacbb feat: implementar HATEOAS en los 9 microservicios que faltaban
9c05bdb feat: agregar validaciones cruzadas via Feign en carrito,
        resenas, pagos y envios
ccaf47d fix: mapear "no encontrado" a HTTP 404 en los 10 microservicios
\end{lstlisting}

\section{Evolución del proyecto: correcciones realizadas durante el semestre}
\label{sec:evolucion}

A diferencia de las evaluaciones parciales anteriores -- que consistieron
exclusivamente en entrega de código y defensa técnica en vivo, sin un informe escrito
--, este documento sí exige explicitar las correcciones realizadas a lo largo del
semestre. El historial de Git del proyecto permite reconstruir esa evolución con
precisión, sin depender de la memoria:

\begin{description}
    \item[Mapeo de errores HTTP.] Hasta antes de la corrección, cualquier recurso no
    encontrado (por ejemplo, un libro con un \texttt{id} inexistente) devolvía
    \texttt{500 Internal Server Error} en lugar de \texttt{404}, porque no existía un
    manejo de excepciones dedicado. Se creó una excepción
    \texttt{RecursoNoEncontradoException} y un \texttt{GlobalExceptionHandler} por
    servicio (reutilizando el que ya existía en catálogo en vez de duplicarlo), lo que
    también reactivó bloques \texttt{catch (FeignException.NotFound)} que hasta
    entonces eran código muerto.
    \item[HATEOAS.] Inicialmente solo el microservicio de catálogo implementaba
    HATEOAS. Se extendió el mismo patrón (\texttt{EntityModel}/\texttt{CollectionModel}
    con enlaces \textit{self}) a los nueve microservicios restantes.
    \item[Validaciones cruzadas vía Feign.] Se agregaron validaciones adicionales entre
    servicios que aún no las tenían: \texttt{carrito} valida contra \texttt{clientes} y
    \texttt{catálogo}; \texttt{reseñas} valida contra \texttt{catálogo} y
    \texttt{clientes}; \texttt{pagos} y \texttt{envíos} validan contra \texttt{pedidos}.
    \item[Cobertura de pruebas.] Se incorporó JaCoCo al POM raíz (antes no existía
    medición real de cobertura) y se agregó una clase de prueba
    \texttt{@WebMvcTest} por servicio, que verifica los códigos HTTP reales devueltos
    por cada endpoint. La cobertura de la lógica de negocio (controladores y
    servicios) es hoy de \textbf{93,6\%} (capítulo~\ref{ch:testing}).
    \item[Relación JPA real entre entidades.] Se agregó al microservicio de envíos una
    segunda entidad, \texttt{HistorialEstado}, en una relación real
    \texttt{@OneToMany}/\texttt{@ManyToOne} con \texttt{Envio}, que registra cada
    cambio de estado del envío. Se optó por este servicio en particular porque la
    relación es puramente aditiva: no modifica el contrato de ningún endpoint
    existente.
    \item[Datos de prueba con Datafaker.] Se incorporó un \texttt{DataSeeder} por
    servicio (\texttt{CommandLineRunner}) que genera datos de ejemplo realistas la
    primera vez que cada tabla está vacía, reutilizando la propia capa de servicio (y
    por lo tanto las validaciones Feign existentes) en lugar de insertar filas
    directamente en el repositorio.
    \item[Validación con Postman.] Se construyó una colección Postman que cubre CRUD y
    comunicación REST entre microservicios para los diez servicios, verificada en vivo
    contra el stack completo levantado en Docker antes de darla por válida
    (capítulos~\ref{ch:crud} y~\ref{ch:rest}).
    \item[Autenticación en clientes.] Se agregó JWT + BCrypt al microservicio de
    clientes, protegiendo únicamente las operaciones de escritura sobre el propio
    perfil (sección~\ref{ch:etica}). Se optó por este alcance acotado en vez de
    proteger todo el microservicio, precisamente porque \texttt{carrito} y
    \texttt{reseñas} dependen de poder leerlo vía Feign sin un token de usuario.
    \item[Reportes de ventas.] Se agregaron reportes de ventas por sucursal y por
    autor en \texttt{pedidos} (sección~\ref{ch:crud}), cerrando parcialmente la
    exclusión de reportes documentada originalmente en la sección~\ref{ch:arquitectura}.
\end{description}

\chapter{Operaciones CRUD con JPA/ORM validadas con Postman}
\label{ch:crud}

\section{Persistencia con Spring Data JPA}

Cada microservicio modela su dominio como una entidad JPA independiente, persistida en
su propia base de datos mediante \texttt{JpaRepository}, sin necesidad de escribir SQL
manual para las operaciones CRUD básicas. El siguiente ejemplo corresponde a la entidad
\texttt{Libro} del microservicio de catálogo:

\begin{lstlisting}[language=Java, caption={Libro.java -- entidad JPA de catálogo}]
@Entity
@Table(name = "libros")
@Data
public class Libro {
    @Id
    @GeneratedValue(strategy = GenerationType.IDENTITY)
    private Long id;

    private String titulo;
    private String autor;
    private Double precio;
}
\end{lstlisting}

\begin{lstlisting}[language=Java, caption={LibroRepository.java -- repositorio}]
public interface LibroRepository extends JpaRepository<Libro, Long> {
}
\end{lstlisting}

La sola extensión de \texttt{JpaRepository<Libro, Long>} otorga las operaciones
\texttt{save}, \texttt{findById}, \texttt{findAll} y \texttt{deleteById} sin
implementación adicional. En los servicios donde la consulta lo requiere se agregan
métodos derivados (por ejemplo, \texttt{findByLibroIdAndSucursal} en inventario, o
\texttt{findByCodigoSeguimiento} en envíos), que Spring Data traduce automáticamente a
JPQL a partir del nombre del método.

Sobre esta base, cada servicio expone las cuatro operaciones CRUD que correspondan a su
dominio a través de endpoints REST. Por ejemplo, el microservicio de clientes expone
las cuatro operaciones completas:

\begin{center}
\begin{tabular}{lll}
\toprule
\textbf{Operación} & \textbf{Método HTTP} & \textbf{Endpoint} \\
\midrule
Crear & POST & \texttt{/api/v1/clientes} \\
Leer (uno) & GET & \texttt{/api/v1/clientes/\{id\}} \\
Leer (todos) & GET & \texttt{/api/v1/clientes} \\
Actualizar & PUT & \texttt{/api/v1/clientes/\{id\}} \\
Eliminar & DELETE & \texttt{/api/v1/clientes/\{id\}} \\
\bottomrule
\end{tabular}
\end{center}

Otros servicios exponen un subconjunto de estas operaciones cuando el dominio no
requiere todas -- por ejemplo, catálogo no expone \texttt{DELETE} porque un libro
referenciado por pedidos o reseñas no debe poder eliminarse sin más.

\section{Validación con la colección Postman}

Para cumplir con el requisito de validar las operaciones CRUD con datos provenientes de
Postman, se construyó una colección (\texttt{postman/BookPoint.postman\_collection.json})
que cubre los diez microservicios, documentada en \texttt{postman/README.md}.

La colección está organizada en once carpetas ordenadas según la cadena real de
dependencias entre servicios (catálogo $\rightarrow$ clientes $\rightarrow$ inventario
$\rightarrow$ carrito $\rightarrow$ pedidos $\rightarrow$ pagos $\rightarrow$ envíos
$\rightarrow$ reseñas $\rightarrow$ notificaciones $\rightarrow$ sucursales), de modo
que puede ejecutarse de principio a fin con \textit{Collection Runner}. Cada request de
creación incluye un script de prueba que captura el \texttt{id} devuelto por la API y
lo guarda en una variable de colección, reutilizada automáticamente por los siguientes
requests -- por ejemplo, el \texttt{id} del libro creado en la carpeta de catálogo se
reutiliza al registrar stock en inventario, agregarlo al carrito y generar un pedido.

\begin{lstlisting}[caption={Script de test de Postman -- captura automática de ID al crear un libro}]
pm.test('Status 201', () => pm.response.to.have.status(201));
const json = pm.response.json();
pm.collectionVariables.set('libroId', json.id);
\end{lstlisting}

Todas las requests de la colección fueron probadas en vivo contra los diez servicios
levantados en Docker antes de considerarlas válidas -- no se trata de una colección
escrita en teoría, sino verificada extremo a extremo, incluyendo los datos generados
automáticamente por Datafaker como insumo (sección~\ref{sec:evolucion}).

\section{Reportes de ventas}

El caso pide que el perfil Jefe de Sucursal pueda "generar reportes de ventas por
categoría, autor o sucursal". Se implementaron los dos primeros en el microservicio de
pedidos, que ya concentra los datos de venta (\texttt{sucursal}, \texttt{libroId},
\texttt{total}):

\begin{itemize}
    \item \texttt{GET /api/v1/pedidos/reportes/por-sucursal}: agregación directa sobre
    la propia base de datos de pedidos (\texttt{GROUP BY sucursal} vía una consulta
    JPQL con proyección de interfaz de Spring Data).
    \item \texttt{GET /api/v1/pedidos/reportes/por-autor}: como el autor no vive en la
    base de datos de pedidos sino en la de catálogo, no es posible resolverlo con un
    \texttt{JOIN} SQL. Se agregan primero las ventas por \texttt{libroId}, y luego se
    enriquece cada grupo con el autor correspondiente vía Feign, reagrupando en memoria
    los libros que comparten autor.
\end{itemize}

El reporte por categoría se dejó fuera: exigiría agregar un campo de categoría o
género al catálogo, un cambio de esquema en otro microservicio no motivado por ninguna
otra necesidad actual del sistema.

Se verificó en vivo, contra el stack completo levantado en Docker con datos generados
por Datafaker, que ambos reportes son consistentes entre sí: la suma total de ventas
agrupada por sucursal coincide exactamente con la suma total agrupada por autor
(\$796.892 en ambos casos sobre el mismo conjunto de pedidos), confirmando que ninguna
venta se pierde ni se duplica al reagrupar.

\chapter{Comunicación RESTful Entre Microservicios}
\label{ch:rest}

\section{Feign Client como mecanismo de comunicación}

La comunicación entre microservicios se implementa de forma declarativa mediante
\textbf{Spring Cloud OpenFeign}, en lugar de construir las llamadas HTTP a mano con
\texttt{RestTemplate} o \texttt{WebClient}. Un \texttt{@FeignClient} se declara como una
interfaz Java; Spring Cloud genera la implementación en tiempo de ejecución y resuelve
la dirección real del servicio destino consultando Eureka -- de modo que ningún
microservicio necesita conocer el host o puerto de otro:

\begin{lstlisting}[language=Java, caption={CatalogoClient.java (inventario) -- Feign client}]
@FeignClient(name = "bookpoint-ms-catalogo")
public interface CatalogoClient {

    @GetMapping("/api/v1/catalogo/{id}")
    EntityModel<LibroRentDTO> obtenerLibroPorId(@PathVariable("id") Long id);
}
\end{lstlisting}

La Tabla~\ref{tab:dependencias-feign} resume las dependencias reales entre
microservicios, todas implementadas mediante este mecanismo.

\begin{table}[h]
\centering
\begin{tabular}{ll}
\toprule
\textbf{Microservicio} & \textbf{Depende vía Feign de} \\
\midrule
Inventario & Catálogo \\
Carrito & Clientes, Catálogo \\
Pedidos & Catálogo, Inventario \\
Pagos & Pedidos \\
Envíos & Pedidos \\
Reseñas & Catálogo, Clientes \\
\bottomrule
\end{tabular}
\caption{Dependencias reales entre microservicios vía Feign}
\label{tab:dependencias-feign}
\end{table}

Los microservicios de sucursales y notificaciones no consumen a ningún otro servicio:
sus dominios son autocontenidos.

\section{Manejo de errores en comunicación remota}

Una llamada Feign puede fallar de dos maneras cualitativamente distintas, y el proyecto
las distingue explícitamente en lugar de tratarlas de forma genérica:

\begin{enumerate}
    \item \textbf{El recurso remoto no existe} (el servicio remoto respondió, pero con
    \texttt{404}): se traduce a una excepción de negocio local
    (\texttt{RecursoNoEncontradoException}), que a su vez el
    \texttt{GlobalExceptionHandler} traduce a un \texttt{404} en la respuesta del
    servicio que hizo la llamada.
    \item \textbf{No se pudo establecer comunicación} (el servicio remoto no responde,
    está caído, o Eureka aún no propagó su registro): se traduce a un error de servidor
    (\texttt{500}), porque no es un problema del dato solicitado sino de la
    infraestructura.
\end{enumerate}

\begin{lstlisting}[language=Java, caption={InventarioServiceImpl.java -- distinción entre 404 remoto y error de red}]
try {
    libroRemoto = catalogoClient.obtenerLibroPorId(inventarioDTO.getLibroId());
} catch (feign.FeignException.NotFound e) {
    throw new RecursoNoEncontradoException(
        "El libro con ID " + inventarioDTO.getLibroId()
        + " no existe en el catalogo.");
} catch (Exception e) {
    throw new RuntimeException(
        "No se pudo conectar con el catalogo. Error de red: "
        + e.getMessage());
}
\end{lstlisting}

Esta distinción no es cosmética: un cliente de la API que recibe un \texttt{404} sabe
que el dato solicitado no existe y puede corregir su petición; uno que recibe un
\texttt{500} sabe que debe reintentar más tarde, porque el problema es de
disponibilidad, no de los datos enviados.

\section{Validación con Postman}

La colección Postman incluye una carpeta dedicada, \textit{"Errores y validaciones
cruzadas (404/400)"}, con casos negativos verificados en vivo contra el stack completo:

\begin{itemize}
    \item Registrar stock de un libro inexistente en inventario $\rightarrow$
    \texttt{404} (validación cruzada contra catálogo).
    \item Crear un pedido con un libro inexistente $\rightarrow$ \texttt{404}
    (validación cruzada contra catálogo).
    \item Crear un pedido con una cantidad mayor al stock disponible
    $\rightarrow$ \texttt{400} ("Stock insuficiente").
    \item Procesar un pago de un pedido inexistente $\rightarrow$ \texttt{404}
    (validación cruzada contra pedidos).
\end{itemize}

Cada uno de estos casos fue ejecutado contra los diez microservicios levantados en
Docker (no simulado) antes de incorporarlo a la colección, incluyendo la verificación
manual de que el mensaje de error y el código HTTP devueltos coincidieran exactamente
con lo documentado en este informe.

\chapter{Pruebas Unitarias}
\label{ch:testing}

\section{Herramientas utilizadas}

Las pruebas automatizadas del proyecto utilizan:

\begin{itemize}
    \item \textbf{JUnit 5} como framework de ejecución.
    \item \textbf{Mockito} para simular (\textit{mock}) repositorios y clientes Feign
    en las pruebas de la capa de servicio, aislándola de la base de datos real y de
    los demás microservicios.
    \item \textbf{@WebMvcTest} + \textbf{@MockitoBean} de Spring Boot Test para las
    pruebas de la capa de controlador, que verifican el código HTTP real devuelto por
    cada endpoint (\texttt{200}, \texttt{201}, \texttt{400}, \texttt{404}) sin levantar
    el contexto completo de Spring ni requerir base de datos.
    \item \textbf{JaCoCo} (\texttt{jacoco-maven-plugin}), integrado en el POM raíz, para
    medir la cobertura real de código ejercitada por las pruebas.
\end{itemize}

\section{Estrategia de pruebas}

Cada servicio cuenta con dos niveles de prueba:

\begin{enumerate}
    \item \textbf{Pruebas de servicio} (\texttt{XServiceImplTest}): verifican la lógica
    de negocio con los repositorios y clientes Feign mockeados, siguiendo el patrón
    \textit{Arrange-Act-Assert}. Cubren tanto el camino feliz como los casos de error
    (recurso no encontrado, argumento inválido).
    \item \textbf{Pruebas de controlador} (\texttt{XControllerTest}): verifican que
    cada endpoint devuelva el código HTTP correcto y la forma esperada de respuesta
    (incluyendo los enlaces HATEOAS), con el servicio mockeado mediante
    \texttt{@MockitoBean}.
\end{enumerate}

El siguiente ejemplo, tomado de las pruebas del microservicio de envíos, ilustra ambos
elementos: la verificación del camino feliz y la del registro correcto en el historial
de estados (la relación JPA agregada en la sección~\ref{sec:evolucion}):

\begin{lstlisting}[language=Java, caption={EnvioServiceImplTest.java -- fragmento}]
@Test
void deberiaCrearEnvioExitosamente() {
    Envio envio = new Envio();
    envio.setPedidoId(1L);
    when(pedidoClient.obtenerPedidoPorId(1L)).thenReturn(new PedidoRentDTO());
    when(envioRepository.save(any(Envio.class)))
        .thenAnswer(invocacion -> invocacion.getArgument(0));

    Envio resultado = envioService.crearEnvio(envio);

    assertEquals("PREPARACION", resultado.getEstado());
    verify(envioRepository, times(1)).save(any(Envio.class));
    verify(historialEstadoRepository, times(1)).save(any(HistorialEstado.class));
}
\end{lstlisting}

\section{Resultados de cobertura}

La ejecución de \texttt{mvn test} sobre los doce módulos del proyecto arroja
\textbf{168 pruebas} en total, todas exitosas. La Tabla~\ref{tab:cobertura} presenta la
cobertura de líneas medida con JaCoCo, calculada de dos formas distintas:

\begin{itemize}
    \item \textbf{Cobertura de lógica de negocio}: solo paquetes
    \texttt{controller}, \texttt{service} y \texttt{mapper} -- el código que realmente
    contiene reglas de negocio y flujo de la aplicación.
    \item \textbf{Cobertura del módulo completo}: incluye además clases de
    configuración (\texttt{WebConfig}, \texttt{OpenApiConfig}), la clase principal de
    arranque de Spring Boot, y los \texttt{DataSeeder} de Datafaker.
\end{itemize}

\begin{table}[h]
\centering
\begin{tabular}{lrr}
\toprule
\textbf{Microservicio} & \textbf{Lógica de negocio} & \textbf{Módulo completo} \\
\midrule
Catálogo & 100,0\% & 71,9\% \\
Clientes & 100,0\% & 70,6\% \\
Inventario & 100,0\% & 53,3\% \\
Notificaciones & 100,0\% & 70,3\% \\
Sucursales & 100,0\% & 70,1\% \\
Pedidos & 95,9\% & 39,6\% \\
Envíos & 93,8\% & 57,8\% \\
Pagos & 88,9\% & 49,5\% \\
Carrito & 81,0\% & 45,4\% \\
Reseñas & 84,3\% & 46,5\% \\
\midrule
\textbf{Total} & \textbf{93,6\%} & \textbf{55,2\%} \\
\bottomrule
\end{tabular}
\caption{Cobertura de líneas por microservicio (JaCoCo)}
\label{tab:cobertura}
\end{table}

La diferencia entre ambas columnas es intencional y se explica completamente: los
\texttt{DataSeeder} agregados para cumplir el requisito de Datafaker (sección
\ref{sec:evolucion}) son clases \texttt{CommandLineRunner} con lógica de reintento
contra Eureka, verificadas mediante ejecución real en Docker contra el stack completo
levantado -- no mediante pruebas unitarias con mocks, ya que su propósito es
precisamente ejercitar la comunicación real entre servicios al arrancar. Contarlas
como "no cubiertas" en un reporte de JaCoCo orientado a pruebas unitarias no refleja
que su comportamiento no haya sido validado, sino que se validó con una herramienta
distinta (verificación en vivo) y no cuenta como cobertura sobre las clases de
configuración (\texttt{WebConfig}, \texttt{OpenApiConfig}) tampoco aporta información
relevante, dado que no contienen lógica de negocio a probar. Por eso se reporta la
cobertura de lógica de negocio como el indicador principal de calidad de las pruebas.

\chapter{Documentación HATEOAS}
\label{ch:hateoas}

\section{Framework elegido: Spring HATEOAS}

Se eligió \textbf{Spring HATEOAS} como framework para implementar el nivel de madurez
de Richardson correspondiente a hipermedia (HATEOAS), por su integración nativa con
Spring Boot: no requiere construir manualmente las URLs de los enlaces, sino que las
deriva directamente de los métodos del propio controlador mediante
\texttt{WebMvcLinkBuilder}, evitando que los enlaces queden desincronizados si cambia
una ruta.

\begin{lstlisting}[language=Java, caption={LibroController.java -- respuesta con enlace HATEOAS}]
@GetMapping("/{id}")
public ResponseEntity<EntityModel<LibroDTO>> obtenerPorId(@PathVariable Long id) {
    Libro libro = libroService.obtenerPorId(id);
    LibroDTO dto = libroMapper.toDTO(libro);

    EntityModel<LibroDTO> model = EntityModel.of(dto);
    model.add(linkTo(methodOn(LibroController.class)
        .obtenerPorId(id)).withSelfRel());
    model.add(linkTo(methodOn(LibroController.class)
        .listar()).withRel("listar"));

    return ResponseEntity.ok(model);
}
\end{lstlisting}

Cada respuesta de un recurso individual incluye al menos un enlace \textit{self}, y en
varios casos un enlace adicional hacia el recurso relacionado más natural (por ejemplo,
de un libro hacia el listado completo, o de una reseña hacia el promedio de
calificaciones del libro correspondiente).

\section{Navegabilidad de la API}

El siguiente ejemplo, tomado de una respuesta real capturada durante la verificación en
Docker, muestra la forma exacta en que Spring HATEOAS serializa un recurso individual
(formato HAL): los campos del recurso se aplanan al nivel superior del JSON, junto al
objeto \texttt{\_links}:

\begin{lstlisting}[language=json, caption={Respuesta real de \texttt{POST /api/v1/catalogo}}]
{
  "_links": {
    "self": { "href": "http://localhost:8081/api/v1/catalogo/3" }
  },
  "id": 3,
  "titulo": "Cien Anos de Soledad",
  "autor": "Gabriel Garcia Marquez",
  "precio": 15990.0
}
\end{lstlisting}

Para una colección, en cambio, el formato HAL cambia de estructura: en lugar de un
arreglo plano, la respuesta envuelve los elementos dentro de \texttt{\_embedded}:

\begin{lstlisting}[language=json, caption={Respuesta real de un listado con HATEOAS}]
{
  "_embedded": {
    "historialEstadoList": [
      { "id": 1, "estado": "PREPARACION", "envioId": 16 },
      { "id": 2, "estado": "EN_CAMINO",  "envioId": 16 }
    ]
  },
  "_links": {
    "self": { "href": "http://localhost:8088/api/v1/envios/16/historial" }
  }
}
\end{lstlisting}

Esta diferencia estructural entre un recurso individual (aplanado) y una colección
(\texttt{\_embedded}) no es cosmética: un consumidor de la API que espere un arreglo
plano de JSON (como un \texttt{List<X>} tipado en un cliente Feign) fallará al
deserializar una respuesta HATEOAS de colección. Este problema se presentó realmente al
extender HATEOAS al microservicio de inventario: el \texttt{InventarioClient} que
consume \texttt{pedidos} esperaba \texttt{List<InventarioRentDTO>}, y tuvo que
actualizarse a \texttt{CollectionModel<EntityModel<InventarioRentDTO>>} para reflejar
la forma real de la respuesta:

\begin{lstlisting}[language=Java, caption={InventarioClient.java (pedidos) -- adaptado a la forma HATEOAS real}]
// inventario expone esta lista envuelta en HATEOAS (CollectionModel), no como
// un arreglo plano -- hay que declarar el tipo tal cual lo devuelve el productor.
@GetMapping("/api/v1/inventario/libro/{libroId}")
CollectionModel<EntityModel<InventarioRentDTO>> obtenerStockPorLibro(
    @PathVariable("libroId") Long libroId);
\end{lstlisting}

Este detalle refuerza por qué la navegabilidad HATEOAS debe considerarse desde el
diseño del contrato entre microservicios, y no agregarse como una capa cosmética al
final del desarrollo.

\chapter{Conclusión}

El sistema desarrollado responde directamente a los tres síntomas identificados en el
caso BookPoint Chile: la lentitud de procesos se aborda mediante microservicios
independientes que escalan y se despliegan sin afectarse entre sí; el stock
desactualizado entre sucursales se resuelve centralizando el inventario en un servicio
único consultado por todos los canales; y las demoras de despacho se abordan con un
microservicio de envíos dedicado, capaz de registrar y exponer el historial completo de
estados de cada envío.

La arquitectura resultante -- diez microservicios independientes, cada uno con su
propia base de datos, comunicándose vía Feign y descubriéndose mediante Eureka detrás de
un API Gateway -- cubre en profundidad el flujo comercial completo del caso: explorar el
catálogo, agregar productos al carrito, comprar, pagar, despachar y reseñar. Quedan
deliberadamente fuera de alcance funciones administrativas de soporte (autenticación y
roles, reportes de venta, promociones, transferencias entre sucursales), priorizando
profundidad y calidad en el flujo principal por sobre cobertura superficial de la
totalidad del caso.

Un aspecto que vale la pena destacar de este proceso es que las correcciones
documentadas en la sección~\ref{sec:evolucion} no fueron anticipadas desde el diseño
inicial, sino descubiertas y corregidas iterativamente a medida que se sometía el
sistema a verificación real -- probando cada flujo contra el stack completo levantado
en Docker, en vez de asumir que el código funcionaba por haber compilado. Esa disciplina
de verificación empírica es, en sí misma, una de las buenas prácticas que este informe
buscó demostrar.

% Espacio para agregar una reflexión personal sobre aprendizajes del proceso, si el
% estudiante quiere incorporar algo aquí antes de la entrega final.


\end{document}
